% =====================================================================
\chapter{IA generativa para 3D y validación de la geometría}
\label{tema:6}
% =====================================================================

Los modelos generativos producen ya objetos 3D a partir de un texto o de una imagen. Para un optimizador geométrico son a la vez una herramienta (generan candidatos o puntos de partida) y un riesgo (sus salidas pueden parecer correctas en una imagen y ser inservibles para fabricar o para un motor). Este tema describe tres familias de generadores con sus artículos originales: la destilación de un modelo de difusión 2D en DreamFusion~\cite{dreamfusion}, los generadores directos de funciones implícitas como Shap-E~\cite{shape} y los generadores autorregresivos de mallas como MeshGPT~\cite{meshgpt}. Después formaliza la validación de la geometría para fabricación y para motores, y la integra en un bucle de generación y verificación. Se apoya en la tarea~4 del curso 16-825 de Carnegie Mellon~\cite{cmu16825} y en el filtro de fabricación aditiva de Langelaar~\cite{langelaar}.

\section{Tres maneras de generar geometría}

\begin{definicion}[Muestreo por destilación de puntuación]\label{def:sds}
Sea $x=g(\theta)$ una imagen renderizada desde una vista aleatoria de una representación 3D con parámetros $\theta$ (en DreamFusion, un campo de radiancia), y sea $\hat\epsilon_\phi(z_t;y,t)$ la predicción del ruido de un modelo de difusión de imágenes preentrenado y congelado, condicionado por el texto $y$, sobre la versión ruidosa $z_t$ de $x$ en el instante $t$. El \emph{muestreo por destilación de puntuación} (SDS) actualiza $\theta$ con el gradiente~\cite[§2.1]{dreamfusion}
\[
\nabla_\theta\mathcal{L}_{\text{SDS}}=\E_{t,\epsilon}\Bigl[w(t)\bigl(\hat\epsilon_\phi(z_t;y,t)-\epsilon\bigr)\frac{\partial x}{\partial\theta}\Bigr],
\]
donde $\epsilon$ es el ruido añadido y $w(t)$ una ponderación.
\end{definicion}

La idea es que el modelo de difusión «sabe» cómo debería parecer una imagen del objeto descrito, y la diferencia entre el ruido que predice y el que realmente se añadió señala en qué dirección corregir la imagen; la regla de la cadena a través del renderizador derivable (\cref{tema:4}) lleva esa corrección a los parámetros 3D. DreamFusion usa el modelo base de $64\times 64$ de Imagen sin modificarlo y un campo de tipo mip-NeRF~360, y encuentra que hace falta una guía sin clasificador muy intensa~\cite{dreamfusion}. En la práctica se mezclan la predicción condicionada y la no condicionada, $\hat\epsilon=\hat\epsilon_u+s(\hat\epsilon_c-\hat\epsilon_u)$, con $s$ grande~\cite{cmu16825}. El método no necesita datos 3D, pero cada objeto es una optimización completa.

Los \emph{generadores directos} aprenden de pares de objetos 3D y textos. Shap-E entrena primero un codificador que transforma cada objeto en los parámetros de una función implícita, que puede renderizarse como malla texturizada o como campo de radiancia, y después un modelo de difusión condicionado sobre esos parámetros; genera objetos en segundos~\cite{shape}. Los \emph{generadores autorregresivos de mallas} tratan la malla como un texto: MeshGPT aprende un vocabulario de \emph{embeddings} geométricos de caras con un codificador de grafos y cuantización vectorial residual, y un transformer de solo decodificador predice el siguiente índice del vocabulario, como un LLM predice el siguiente token (módulos~12 y~15)~\cite{meshgpt}. El resultado son mallas compactas, parecidas a las que haría un artista, a diferencia de las mallas densas que se obtienen extrayendo la isosuperficie de un campo; a cambio, el tamaño está limitado por el contexto (los autores entrenan con mallas de menos de 800 caras).

Ninguna de las tres familias garantiza que el resultado sea una malla cerrada, orientada, sin autointersecciones, con la escala correcta y fabricable. Esas propiedades hay que comprobarlas.

\section{Validación de la geometría}

\begin{definicion}[Comprobaciones topológicas]\label{def:validacion}
Sea $(V,F)$ una malla triangular. Se cuentan las \emph{aristas dirigidas} $(a,b)$ de todas las caras (cada cara $(a,b,c)$ aporta $(a,b)$, $(b,c)$ y $(c,a)$) y las aristas no dirigidas $\{a,b\}$. La malla
\begin{enumerate}
  \item tiene \emph{borde} si alguna arista no dirigida pertenece a una sola cara;
  \item es \emph{no variedad} si alguna arista pertenece a tres o más caras;
  \item está \emph{mal orientada} si alguna arista dirigida aparece más de una vez.
\end{enumerate}
Una malla cerrada, variedad y orientada de forma coherente es \emph{estanca}.
\end{definicion}

\begin{proposicion}[Una cara invertida]\label{prop:cara-girada}
Sea una malla estanca e invértase la orientación de una única cara $f$ cuyas tres aristas son compartidas con caras distintas de $f$. Entonces la malla sigue cerrada y sin aristas no variedad, exactamente tres aristas dirigidas aparecen dos veces, y el volumen de la \cref{prop:volumen} disminuye en $\tfrac13\det(v_a,v_b,v_c)$, el doble de la contribución original de $f$.
\end{proposicion}

\begin{proof}
Las aristas no dirigidas no cambian, luego el recuento por arista sigue siendo dos: la malla sigue cerrada y es variedad. Antes, cada arista $\{a,b\}$ de $f$ se recorría como $(a,b)$ en $f$ y como $(b,a)$ en su vecina; tras invertir $f$, ambas la recorren como $(b,a)$, que aparece dos veces; esto ocurre para las tres aristas de $f$ y para ninguna otra. Finalmente, invertir la cara intercambia dos columnas del determinante, que cambia de signo: la suma pierde $\det/6$ y gana $-\det/6$, una variación de $-\det/3$.
\end{proof}

Por eso una cara invertida es un defecto silencioso: no abre agujeros y una comprobación de estanqueidad basada solo en recuentos de aristas no dirigidas no la detecta, pero estropea el volumen, la iluminación (la normal apunta hacia dentro) y la clasificación interior/exterior que usan los programas de laminado.

Para \emph{fabricación aditiva} se añaden comprobaciones geométricas:
\begin{itemize}[leftmargin=1.4em]
  \item \emph{Voladizos}: con dirección de fabricación $b$ unitaria, una cara de normal $n$ que mira hacia abajo necesita soporte si forma con la horizontal un ángulo menor que un umbral, habitualmente $45^\circ$; equivale a $n\cdot b<-\cos 45^\circ\approx-0{,}707$. Las caras apoyadas en la bandeja se excluyen. El filtro de Langelaar incorpora una restricción de este tipo dentro de la optimización de topología~\cite{langelaar}.
  \item \emph{Espesor mínimo}: con la SDF de la pieza (\cref{tema:1}), el espesor local es aproximadamente el doble de la distancia al borde en el eje medial; en la práctica se lanzan rayos desde la superficie hacia dentro y se mide la distancia a la cara opuesta.
  \item \emph{Volumen de impresión, unidades y escala}: la caja envolvente debe caber en la máquina y estar en milímetros.
  \item \emph{Cavidades cerradas}: en procesos de lecho de polvo o resina atrapan material sin consolidar.
\end{itemize}
Para un \emph{motor de videojuegos}: número de triángulos y de materiales dentro del presupuesto (\cref{tema:5}), coordenadas UV presentes y sin solapes, normales y tangentes coherentes, escala y punto de pivote convenidos, niveles de detalle generados y nombres según la convención del proyecto.

\section{Generar y verificar}

El esquema robusto es un bucle: el generador propone, el validador comprueba, un reparador corrige lo automático (soldar vértices duplicados, reorientar caras, cerrar agujeros pequeños) y el resultado se vuelve a validar. Si cada candidato pasa la validación con probabilidad $p$, generar varios y quedarse con el primero válido es eficaz (\cref{fig:bucle}).

\begin{figure}[h]
\centering
\begin{tikzpicture}[every node/.style={font=\small},
  caja/.style={draw, rounded corners=2pt, minimum height=0.85cm, align=center, inner sep=4pt},
  flecha/.style={-{Stealth[length=2.2mm]}, thick}]
\node[caja] (gen) at (0,0) {generador\\ (texto o imagen)};
\node[caja] (val) at (4.0,0) {validador de\\ selección};
\node[caja] (ind) at (8.0,0) {comprobaciones\\ independientes};
\node[caja] (ok) at (11.6,0) {pieza o\\ recurso};
\node[caja] (rep) at (4.0,-2.2) {reparador\\ automático};
\draw[flecha] (gen) -- (val);
\draw[flecha] (val) -- node[above] {pasa} (ind);
\draw[flecha] (ind) -- node[above] {pasa} (ok);
\draw[flecha] ([xshift=-0.35cm]val.south) -- node[left] {falla} ([xshift=-0.35cm]rep.north);
\draw[flecha] ([xshift=0.35cm]rep.north) -- node[right] {revalidar} ([xshift=0.35cm]val.south);
\draw[flecha, dashed] (ind.south) |- node[pos=0.75, above, align=center] {fallo: ampliar\\ la selección} (rep.east);
\draw[flecha, dotted] (rep.west) -| node[pos=0.25, above, align=center] {irreparable:\\ otro candidato} (gen.south);
\end{tikzpicture}
\caption{Bucle de generación y verificación. Las comprobaciones independientes no se usan para seleccionar ni para ajustar el generador.}
\label{fig:bucle}
\end{figure}

\begin{proposicion}[Lo mejor de $N$]\label{prop:mejorN}
Si $N$ candidatos superan la validación de forma independiente con probabilidad $p\in(0,1)$, la probabilidad de obtener al menos uno válido es $1-(1-p)^N$, y para alcanzar una probabilidad $1-\delta$ basta con
\[
N\ge\frac{\ln\delta}{\ln(1-p)} .
\]
\end{proposicion}

\begin{proof}
Por independencia (0.8, §4.4), la probabilidad de que fallen todos es $(1-p)^N$. La condición $(1-p)^N\le\delta$ equivale, tomando logaritmos y dividiendo por $\ln(1-p)<0$, que invierte la desigualdad, a $N\ge\ln\delta/\ln(1-p)$.
\end{proof}

El validador se convierte así en una función de recompensa verificable, análoga a las recompensas verificables con que se ajustan los LLM (módulo~21, §6.2). Y hereda su riesgo: si el generador se ajusta para pasar el validador, aprenderá a satisfacer las comprobaciones sin satisfacer la intención (por ejemplo, mallas estancas que se autointersecan si no se comprueban las autointersecciones). Es la ley de Goodhart que el módulo~21 (§6.4) estudia como sobreoptimización de la recompensa. La defensa es la misma: comprobaciones independientes que no se usan para ajustar el generador y revisión humana por muestreo.

\section{Caso de estudio: validación de mallas generadas}

\begin{ejemplo}[Un cubo con defectos y un voladizo]\label{ej:validacion}
\emph{Contexto.} Un generador produce una pieza que debería ser un cubo de 40 mm de lado centrado en el origen, con 8 vértices y 12 triángulos. Se validan tres salidas: la correcta, una a la que le falta un triángulo y una con un triángulo invertido.

\emph{Comprobaciones} (\cref{def:validacion} y \cref{prop:volumen}):

\begin{center}\small
\begin{tabular}{@{}lrrrrrr@{}}
\toprule
Salida & $\abs{E}$ & borde & no variedad & dirigidas repetidas & $\chi$ & volumen (mm³) \\
\midrule
correcta & 18 & 0 & 0 & 0 & 2 & 64\,000 \\
falta un triángulo & 18 & 3 & 0 & 0 & 1 & 58\,667 \\
un triángulo invertido & 18 & 0 & 0 & 3 & 2 & 53\,333 \\
\bottomrule
\end{tabular}
\end{center}

\emph{Lectura.} La malla incompleta tiene 3 aristas de borde y $\chi=1$; su «volumen» ($64\,000-5333$) no significa nada porque la superficie no encierra una región. La malla con un triángulo invertido pasa las pruebas de borde y de variedad, pero tiene 3 aristas dirigidas repetidas y su volumen es $\tfrac56$ del real: cada uno de los 12 triángulos aporta $a^3/12=5333$ mm³ y, por la \cref{prop:cara-girada}, invertir uno resta el doble. Una validación que solo contara aristas no dirigidas la daría por buena.

\emph{Voladizos.} Para imprimir el cubo en $+z$ inclinado $30^\circ$ alrededor del eje $x$, las normales de las caras que miran hacia abajo tienen componente vertical $-\cos 30^\circ=-0{,}866$ (la antigua base, que ya no apoya en la bandeja) y $-\sin 30^\circ=-0{,}5$. Solo la primera cumple $n_z<-0{,}707$: necesita soporte $1600$ mm² de $9600$, un $16{,}7$~\%. La segunda forma $60^\circ$ con la horizontal y se autosoporta. Apoyar el cubo sobre una cara elimina los soportes.

\emph{Política de generación.} Si el generador produce piezas válidas con probabilidad $p=0{,}3$ e independencia entre muestras, la \cref{prop:mejorN} con $\delta=0{,}05$ da $N\ge\ln 0{,}05/\ln 0{,}7=8{,}4$: generando 9 candidatos se obtiene al menos uno válido con probabilidad $0{,}960$ (con 8 sería $0{,}942$).
\end{ejemplo}

\begin{table}[h]
\centering\small
\begin{tabularx}{\textwidth}{@{}l>{\raggedright\arraybackslash}X>{\raggedright\arraybackslash}p{2.6cm}>{\raggedright\arraybackslash}X@{}}
\toprule
Herramienta & Para qué sirve & Licencia o tipo & Cuándo elegirla \\
\midrule
trimesh & estanqueidad, orientación, volumen, reparaciones simples & MIT & validador ligero en Python \\
PyMeshLab & filtros de reparación (no variedad, agujeros) & GPL-3.0 & reparación por lotes \\
Open3D & métricas geométricas y distancias entre mallas y nubes & MIT & comparar con una referencia \\
Blender (\texttt{bpy}) & comprobaciones y correcciones dentro del pipeline del artista & GPL-3.0 & validar recursos antes de exportarlos al motor \\
Shap-E & generación de funciones implícitas a partir de texto o imagen & MIT (pesos publicados) & prototipos y generación local reproducible \\
\bottomrule
\end{tabularx}
\caption{Tecnologías del tema~6.}
\label{tab:tec6}
\end{table}

\begin{observacion}[Errores frecuentes en la práctica]\label{obs:errores6}
(i) \emph{Validar por la imagen}: un render convincente no garantiza una malla estanca, orientada y sin autointersecciones. (ii) \emph{Escala y unidades}: los generadores suelen producir objetos normalizados; sin una escala explícita, una pieza de 40 mm puede llegar a la impresora como de 40 m o de 0,04 mm. (iii) \emph{Validar con los mismos casos con que se ajustó el generador}: los textos o imágenes de prueba deben ser distintos de los usados para ajustar o elegir el generador. (iv) \emph{Goodhart}: un generador seleccionado o ajustado contra el validador aprende sus huecos; se reservan comprobaciones independientes. (v) \emph{Propiedad intelectual y licencias}: los pesos y los datos de entrenamiento de los generadores tienen licencias propias que limitan el uso comercial de los resultados en algunos casos. (vi) \emph{Reparaciones que cambian la pieza}: cerrar agujeros o soldar vértices puede alterar dimensiones funcionales; tras reparar se vuelve a medir.
\end{observacion}

\begin{lecturas}
\item B.~Poole \emph{et al.}~\cite{dreamfusion}: §2--3 (muestreo por destilación de puntuación y el modelo de NeRF).
\item H.~Jun y A.~Nichol~\cite{shape}: generación de funciones implícitas condicionadas.
\item Y.~Siddiqui \emph{et al.}~\cite{meshgpt}: generación de mallas con transformers de solo decodificador.
\item CMU 16-825~\cite{cmu16825}, tarea~4, §2 «Diffusion-guided Optimization».
\item M.~Langelaar~\cite{langelaar}: restricción de voladizo en la optimización de topología.
\end{lecturas}

\section*{Ejercicios del tema 6}
\addcontentsline{toc}{section}{Ejercicios del tema 6}

\begin{ejercicio}{1}{eje:6-checks}
Enumere cuatro comprobaciones para una pieza que se va a imprimir y cuatro para un recurso que se va a cargar en un motor. ¿Cuáles son comunes?
\end{ejercicio}

\begin{ejercicio}{2}{eje:6-sds}
En un ejemplo de juguete con dos píxeles, el ruido añadido es $\epsilon=(0{,}3;\ -0{,}2)$, la predicción condicionada es $\hat\epsilon_c=(0{,}5;\ 0{,}1)$ y la no condicionada $\hat\epsilon_u=(0{,}2;\ -0{,}1)$. Con guía $s=7{,}5$, $w(t)=1$ y $x=\theta$ (se optimiza directamente la imagen), calcule $\hat\epsilon$, el gradiente SDS y el nuevo $\theta$ tras un paso de descenso con tasa $0{,}01$ desde $\theta=(0{,}5;\ 0{,}5)$.
\fuente{adaptado de CMU 16-825, «Learning for 3D Vision» (S.~Tulsiani), tarea~4, §2.1 «SDS Loss + Image Optimization»~\cite{cmu16825}.}
\end{ejercicio}

\begin{ejercicio}{2}{eje:6-aristas}
Una malla tiene las caras $(0,1,2)$, $(0,2,3)$, $(0,3,1)$, $(1,3,2)$ y $(1,2,4)$. Clasifique sus aristas (de borde, interiores, no variedad), calcule $\chi$ y diga qué defectos tiene.
\end{ejercicio}

\begin{ejercicio}{2}{eje:6-voladizo}
Una pieza se imprime en $+z$. Sus caras tienen áreas y componentes $n_z$: 400 mm² con $n_z=-1$ (apoyada en la bandeja), 300 mm² con $n_z=-0{,}8$, 200 mm² con $n_z=-0{,}6$, 500 mm² con $n_z=0$ y 600 mm² con $n_z=0{,}5$. Con umbral de $45^\circ$, ¿qué fracción del área total necesita soporte?
\end{ejercicio}

\begin{ejercicio}{2}{eje:6-mejorN}
Un generador produce piezas válidas con probabilidad $0{,}2$. Cada generación tarda 30 s y cada validación 2 s. ¿Cuántos candidatos hay que generar para tener al menos uno válido con probabilidad $0{,}9$, y cuánto tarda el lote? ¿Cuál es el número esperado de candidatos hasta el primer válido si se generan de uno en uno?
\end{ejercicio}

\begin{ejercicio}{2}{eje:6-mc}
El SDF de un campo neuronal se evalúa en una rejilla de $128^3$ nodos sobre $[-1{,}5;\ 1{,}5]^3$ para extraer la malla. (a) Calcule el paso de la rejilla y el número de evaluaciones, y cuántos lotes de 65\,536 puntos hacen falta. (b) Si una unidad equivale a 50 mm, ¿se resolverá una pared de 0,8 mm?
\fuente{adaptado de Stanford CS348n (L.~Guibas, primavera de 2023), tarea~2, problema~1(c) «Marching cubes»~\cite{cs348nhw2}.}
\end{ejercicio}

\begin{ejercicio}{3}{eje:6-genero}
(a) Para una malla triangular cerrada y conexa de género $g$, demuestre que $\abs{V}-\tfrac12\abs{F}=2-2g$. (b) Un toro se discretiza con una rejilla de $16\times 8$ vértices con identificación periódica en ambas direcciones, dividiendo cada cuadrilátero en dos triángulos. Calcule $\abs{V}$, $\abs{E}$, $\abs{F}$, $\chi$ y $g$. (c) Proponga cómo usar $g$ para validar la salida de un generador que debería producir una arandela.
\end{ejercicio}

\begin{ejercicio}{3}{eje:6-goodhart}
Diseñe razonadamente un bucle de generación y validación para producir soportes imprimibles a partir de descripciones de texto: indique qué comprobaciones usa para seleccionar candidatos, cuáles reserva como comprobaciones independientes, qué reparaciones automáticas admite y cómo detectaría que el generador ha empezado a «engañar» al validador.
\end{ejercicio}
